This report studies defect-aware acceleration for a practical 3D
semiconductor X-ray inspection workflow composed of \emph{wp5-seg} for
segmentation-model training and compression, and \emph{intelliscan}
for end-to-end raw-scan inference, metrology, reporting, and deployment.
The core problem is that patch-level speed or model-level Dice does not
guarantee safe whole-pipeline behavior. In this application, small
segmentation changes can trigger large metrology changes and binary
defect-flag flips.

To address this problem, the report adopts a \textbf{Two-Level
Deployment-Safety Gate}. At Gate 1, methods are screened on the full
174-case \emph{wp5-seg} evaluation set. At Gate 2, retained
methods are re-evaluated inside \emph{intelliscan} with emphasis on
class-sensitive drift, metrology deltas, and defect flips. Using this
rule, the report identifies both retained positive results and validated
negative results.

On the engineering side, combined segmentation plus metrology, report
and GLTF artifact cleanup, exact quantized in-memory detection,
4-worker CPU image preparation, and intra-view batch prefetch
were all validated. On the current local eight-sample raw-scan set, the
final safe default path reduces mean total runtime by 45.13\%
and pooled total runtime by 44.20\% relative to the older
file-based detection baseline while preserving exact outputs.

On the model side, structured pruning and class-aware knowledge
distillation established a first Pareto front. Pruning ratio
0.375 is the current quality-oriented student, while
0.5 is the strongest balanced point. However, deployment
experiments show that even 0.8895 full-case average Dice can
still cause unacceptable downstream defect flips in the whole pipeline.
Additional ablations demonstrate that context shrinkage, not the real
fast path itself, is the dominant cause of adaptive-margin drift, and
that larger direct padded ROI shells are also unsafe under the current
model.

The main contribution of the report is therefore a defect-aware
co-optimization methodology for the full 3D inspection stack, rather
than a single faster network.
